% \nocite{*}
% \printbibliography
\begin{thebibliography}{9}
\bibitem{grantmestrength_opengl}
  Microsoft Learn. (s.f.). \textit{Control gráfico OpenGL - Win32 apps}. Recuperado de \url{https://learn.microsoft.com/es-es/windows/win32/opengl/introduction-to-opengl}

\bibitem{microsoft_rgba}
Microsoft Learn. (s.f.). \textit{Modo RGBA y administración de paletas de Windows - Win32 apps}. Recuperado de \url{https://learn.microsoft.com/es-es/windows/win32/opengl/rgba-mode-and-windows-palette-management}

\bibitem{hearn_baker}
Hearn, D., \& Baker, M. P. (2004). \textit{Computer graphics with OpenGL} (3.\textsuperscript{a} ed.). Pearson Prentice Hall.

\bibitem{marschner_shirley}
Marschner, S., \& Shirley, P. (2021). \textit{Fundamentals of computer graphics} (5.\textsuperscript{a} ed.). CRC Press.

\bibitem{shreiner_opengl}
Shreiner, D., Sellers, G., Kessenich, J., \& Licea-Kane, B. (2013). \textit{OpenGL programming guide: The official guide to learning OpenGL, version 4.3} (8.\textsuperscript{a} ed.). Addison-Wesley.

\bibitem{learnopengl_hello}
de Vries, J. (s.f.). \textit{Hello triangle}. LearnOpenGL. Recuperado de \url{https://learnopengl.com/Getting-started/Hello-Triangle}

\bibitem{khronos_polygonmode}
Khronos Group. (s.f.). \textit{glPolygonMode - OpenGL 4 Reference Pages}. Recuperado de \url{https://registry.khronos.org/OpenGL-Refpages/gl4/html/glPolygonMode.xhtml}

\bibitem{glm_manual}
G-Truc Creation. (s.f.). \textit{OpenGL Mathematics (GLM) manual}. Recuperado de \url{https://glm.g-truc.net/0.9.9/index.html}

\bibitem{glfw_window}
GLFW. (s.f.). \textit{Window guide: Buffer swapping}. Recuperado de \url{https://www.glfw.org/docs/latest/window_guide.html}

\bibitem{cppreference_init}
cppreference.com. (s.f.). \textit{Default initialization}. Recuperado de \url{https://en.cppreference.com/w/cpp/language/default_initialization}

\bibitem{teodoro_manual03}
Teodoro Vite, S. (s.f.). \textit{Manual de prácticas. Práctica 03: Modelado geométrico} [Manual de laboratorio]. Facultad de Ingeniería, Universidad Nacional Autónoma de México.

\bibitem{anthropic2026}
Anthropic. (2026). \textit{Claude} (Opus 5) [Modelo de lenguaje de gran escala]. \url{https://claude.ai}
\textit{Se utilizó Claude como apoyo para analizar el comportamiento de la
clase Circle del proyecto base y en la revisión de redacción del presente
reporte. El código base fue proporcionado por el profesor; las modificaciones
al mismo, las capturas y los resultados son de autoría propia.}

\end{thebibliography}
